\documentclass[manuscript]{acmart}

% For anonymous review, use the options required by the current TOSN call:
% \documentclass[manuscript,review,anonymous]{acmart}

\usepackage{booktabs}
\usepackage{graphicx}
\usepackage{xspace}

\newcommand{\smartwilds}{\textsc{SmartWilds}\xspace}

\begin{document}

\title{\smartwilds: A Systems Dataset and Metadata Foundation for
Cross-Modal Ecological Sensing}

% Alternative titles:
% \title{\smartwilds: Traces, Benchmarks, and Metadata for
% Cross-Modal Ecological Sensing}
%
% \title{From Environmental Sensors to Coordinated Sensing Systems:
% The \smartwilds Dataset and Benchmark}

% Claim: environmental monitoring established the infrastructure → ecological monitoring specialized it around living systems → smart sensors make coordination possible → SmartWilds exposes the systems and data foundations that coordination requires.

\begin{abstract}
Recent advances in environmental sensing have made it possible to
observe ecosystems using increasingly capable networks of camera
traps, acoustic recorders, drones, animal-borne sensors, and edge
computers. These technologies build on decades of environmental
sensing research, but their combination creates a new systems
opportunity: coordinating heterogeneous sensors so that observations
from one modality can inform when, where, and how other sensors
collect and process data.

Studying such systems requires datasets that preserve more than sensor
observations and ecological labels. Systems researchers also need the
timing, configuration, spatial coverage, provenance, missingness, and
operational context through which observations were produced. Existing
ecological and computer-vision datasets rarely retain this information.

We present \smartwilds, a synchronized, cross-modal dataset collected
during a field deployment at The Wilds conservation center. The release
includes drone imagery, camera-trap observations, bioacoustic
recordings, location and environmental context, deployment records, and
systems-oriented traces. We characterize the resulting workloads,
including their burstiness, temporal and spatial overlap, data demand,
and opportunities for cross-modal coordination. We further introduce a
benchmark suite for trace-driven systems evaluation and a comprehensive
dataset card that documents the metadata required to interpret, replay,
and extend cross-modal field deployments. Together, these contributions
provide an empirical foundation for designing adaptive ecological
sensing systems and a starting point for interoperable dataset
standards that preserve both ecological meaning and systems provenance.
\end{abstract}

\keywords{environmental sensing, ecological monitoring, sensor
networks, edge computing, multimodal sensing, cross-modal coordination,
datasets, benchmarks, metadata, record and replay}

\maketitle

% ============================================================
% ============================================================
\section{Introduction}
\label{sec:introduction}

% Environmental sensing has laid the groundwork for distributed,
% long-term observation under field constraints.
%
% Recent advances in low-cost sensors, embedded computing, and AI now
% make it practical to combine stationary, mobile, and
% animal-associated sensing modalities.
%
% These modalities provide complementary views, but collecting them
% together does not automatically produce an integrated systems
% dataset. They differ in spatial support, temporal cadence, timestamp
% semantics, data format, sampling policy, uncertainty, and provenance.
%
% Core question:
% What systems and data infrastructure is required to construct a
% reusable mesoscale sensing dataset from heterogeneous ecological
% sensors?
%
% Introduce SmartWilds as:
% 1. a coordinated multimodal field campaign;
% 2. an implemented cross-modal data infrastructure;
% 3. a complete, systems-ready dataset release;
% 4. a source of generalizable design requirements.

\subsection{Research Questions}

\begin{description}
    \item[RQ1:] What information must be preserved to cross-reference
    observations from heterogeneous ecological sensors?

    \item[RQ2:] How can a shared data model integrate these
    observations without discarding their modality-specific spatial,
    temporal, and semantic properties?

    \item[RQ3:] Can the resulting infrastructure support provenance
    reconstruction, cross-modal queries, workload characterization,
    and trace-driven systems evaluation?

    \item[RQ4:] What lessons from SmartWilds generalize to future
    mesoscale environmental sensing deployments?
\end{description}

\subsection{Contributions}

\begin{itemize}
    \item \textbf{A systems framing of cross-modal ecological sensing.}
    We situate cross-modal ecological sensing within the broader
    environmental sensing literature and identify the additional
    infrastructure required to relate complementary biological
    observations across modalities and scales.

    \item \textbf{A cross-modal data infrastructure.}
    We design an entity and relationship model, identifier scheme,
    provenance representation, and processing pipeline for preserving
    and cross-referencing heterogeneous field observations.

    \item \textbf{The complete SmartWilds systems dataset.}
    We release multimodal observations, collection protocols, device
    configurations, deployment records, annotations, systems traces,
    provenance, and known limitations from a coordinated field
    campaign.

    \item \textbf{An empirical systems evaluation.}
    We evaluate the infrastructure through metadata completeness,
    provenance reconstruction, representative cross-modal queries,
    workload characterization, and generation of replayable traces.

    \item \textbf{Documentation and design guidance.}
    We provide a comprehensive dataset card and derive recommendations
    for constructing future interoperable mesoscale sensing datasets.
\end{itemize}

\newpage
% ============================================================
\section{From Environmental Sensing to Cross-Modal Ecological Sensing}
\label{sec:context}

\subsection{Environmental Sensing Laid the Groundwork}

% Distributed sensing, wireless sensor networks, long-term
% environmental observatories, mobile sensing, remote instrumentation,
% edge computing, data collection under field constraints.

\subsection{Ecological Sensing Builds on This Foundation}

% Organisms, behavior, communities, interactions, ecosystem processes.
% Mobile subjects, rare events, observer effects, semantic uncertainty.

\subsection{Complementary Sensors for Mesoscale Observation}

% Camera traps, acoustics, drones, GPS, environmental measurements.
% Complementarity across coverage, cadence, mobility, and semantics.
%
% Do not present the ecological results from the MEE paper.
% The point here is the systems consequence of complementarity.

\subsection{Why Multimodal Collection Is Not Yet Cross-Modal Sensing}

% Multiple datasets can coexist without being interoperable.
% Explain mismatched:
% - timestamps;
% - spatial footprints;
% - sampling effort;
% - identifiers;
% - file structures;
% - uncertainty;
% - transformations;
% - provenance.

\subsection{Requirements for a Systems-Ready Dataset}

% Joint discovery and querying.
% Modality-specific semantics.
% Provenance and uncertainty.
% Missingness and failures.
% Human and computational processes.
% Trace construction and future coordination studies.


\newpage
% ============================================================
\section{Related Work}
\label{sec:related}

\subsection{Environmental Sensor Networks and Field Deployments}

\subsection{Multimodal Ecological Sensing and Datasets}

\subsection{Sensor Data Management and Semantic Querying}

\subsection{Metadata, Provenance, and Ecological Standards}

% Darwin Core, FAIR, FAIR4AI, FAIR^2 Drones,
% dataset cards, modality-specific standards.

\subsection{Systems Traces, Replay, and Digital Twins}

% Position replay as an established systems methodology.
% SmartWilds provides field-derived workloads for its application to
% emerging cross-modal sensing systems.

\newpage
% ============================================================
\section{SmartWilds Study Design and Data Collection}
\label{sec:collection}

% This section establishes what produced the dataset.
% Keep ecological results out.

\subsection{Study Site and Observation Objectives}

\subsection{Spatial and Temporal Sampling Design}

\subsection{Sensing Modalities}

\subsubsection{Drone Imagery and Telemetry}
\subsubsection{Camera-Trap Images and Video}
\subsubsection{Bioacoustic Recordings}
\subsubsection{GPS and Location Data}
\subsubsection{Weather and Environmental Context}

\subsection{Device Configuration and Field Protocols}

\subsection{Human Roles and Operational Workflow}

\subsection{Clock Synchronization and Timestamp Semantics}

\subsection{Spatial Referencing and Sensor Coverage}

\subsection{Ethics, Permissions, and Sensitive Information}

% ============================================================
\newpage
\section{Cross-Modal Data Infrastructure}
\label{sec:infrastructure}


% This should be one of the paper's central technical sections.

\subsection{Design Requirements}

% Heterogeneity, uncertainty, traceability, extensibility,
% machine-readability, compatibility with ecological standards.

\subsection{Architecture and Data Flow}

% From devices and field records through ingestion, validation,
% processing, indexing, manifests, and release.

\begin{figure}[t]
    \centering
    % \includegraphics[width=\linewidth]{figures/data-infrastructure.pdf}
    \caption{SmartWilds data infrastructure from heterogeneous field
    collection through ingestion, provenance-preserving processing,
    cross-modal indexing, and dataset release.}
    \label{fig:data-infrastructure}
\end{figure}

\subsection{Core Entity Model}

% Site, campaign, deployment, device, session, observation, asset,
% annotation, ecological entity, relationship, transformation,
% system event, human action.

\subsection{Identifiers and Cross-Modal Relationships}

% Stable IDs.
% Relationships do not presume exact ecological co-occurrence.
% Temporal relation, spatial relation, common context,
% derived-from, annotated-by, processed-by, and triggered-by.

\subsection{Temporal and Spatial Representation}

% Timestamp meaning, interval, clock source, uncertainty.
% Point, trajectory, field of view, acoustic region, aerial footprint,
% coordinate reference system, spatial uncertainty.

\subsection{Provenance Model}

% Raw -> transformed -> inferred -> annotated -> associated.
% Software, model, parameters, person, time, and source assets.

\subsection{QA/QC: Missingness, Failure, and Sampling Effort}

% Distinguish:
% - no ecological observation;
% - sensor inactive;
% - sensor active but no trigger;
% - lost/corrupted data;
% - metadata unavailable;
% - observation excluded during quality control.

\subsection{Implementation}

% Storage formats, manifests, indexes, APIs/loaders, validation tools,
% checksums, versioning, and repository organization.
%
% This subsection is essential for TOSN: describe what was actually
% built, not only what a future schema should contain.

\newpage
% ============================================================
\section{The SmartWilds Dataset Release}
\label{sec:dataset}

\subsection{Dataset Composition}

% Counts, duration, data volume, modalities, sessions, formats.

\subsection{Released Data Products}

% Raw observations, processed products, annotations, deployment
% metadata, field logs, configurations, traces, documentation, code.

\subsection{Machine-Readable Metadata}

% Example record and data dictionary.
% Place the complete field dictionary in the artifact/supplement.

\subsection{Dataset Organization and Versioning}

\subsection{Access, Licensing, Maintenance, and Citation}

\subsection{Known Limitations}

% State explicitly that ecological analysis and interpretation are
% reported separately.

\newpage
% ============================================================
\section{Systems-Oriented Dataset Characterization}
\label{sec:characterization}

% Characterize the workload without duplicating the ecological study.

\subsection{Data Volume and Composition}

\subsection{Arrival Patterns and Burstiness}

\subsection{Sensor Operating Schedules}

\subsection{Spatial and Temporal Coverage}

% Coverage as a property of the sensing system.
% Do not reproduce ecological co-occurrence findings.

\subsection{Missingness and Operational Failures}

\subsection{Storage, Processing, and Data-Movement Requirements}

\subsection{Human Effort}
% Setup, collection, transfer, validation, annotation, and repair.

\newpage
% ============================================================
\section{Evaluation}
\label{sec:evaluation}

% Organize explicitly around the research questions.

\subsection{Metadata Coverage and Validation}

% Questions:
% - What proportion of released assets resolve to a device, session,
%   time reference, spatial reference, and provenance record?
% - Which fields were recorded, reconstructed, inferred, or missing?
% - Does the validator detect broken references and invalid records?

\subsection{Cross-Modal Query Support}

% Representative systems queries:
% - Retrieve all active sensors and observations for a site/time range.
% - Retrieve data products spatially related to a drone trajectory.
% - Retrieve source observations underlying a derived record.
% - Distinguish absent observations from inactive or failed sensors.
%
% Evaluate correctness and, where useful, latency/scalability.

\subsection{Provenance Reconstruction}

% Select representative derived products.
% Demonstrate reconstruction back to raw assets, device configuration,
% processing software/model, parameters, and human annotations.

\subsection{Trace Generation and Replay}

% Show that machine-readable dataset records can generate ordered,
% modality-specific workloads.
% Replay one representative interval or workload.
% Report reproducibility and resource metrics.

\subsection{Comparison with Existing Metadata Standards}

% Map required fields to existing standards.
% Identify:
% - reused fields;
% - specialized fields;
% - systems extensions.
%
% Compare expressiveness, not merely the number of fields.

\subsection{Summary of Findings}

% Answer RQ1--RQ4 directly.

\newpage
% ============================================================
\section{Dataset Documentation and Standards Alignment}
\label{sec:documentation}

% This section summarizes the dataset-card contribution.
% The complete dataset card should accompany the release.

\subsection{Dataset Card Structure}

\subsection{Alignment with Darwin Core, FAIR, and FAIR$^2$ Drones}

\subsection{Systems Metadata Not Captured Consistently Elsewhere}

\subsection{Human- and Machine-Readable Documentation}
% Avoid claiming a finished universal standard.
% Present this as a requirements analysis, reference implementation,
% and foundation for further community development.

\newpage
% ============================================================
\section{Lessons for Future Proximal (Mesoscale?) Sensing Datasets}
\label{sec:lessons}

\subsection{Plan Identifiers and Relationships Before Deployment}

\subsection{Record Timestamp Meaning and Uncertainty}

\subsection{Represent Coverage and Sampling Effort}

\subsection{Preserve Modality-Specific Semantics}

\subsection{Make Missingness and Failure First-Class Data}

\subsection{Track Human and Computational Provenance}

\subsection{Design Data Infrastructure as Part of the Sensing System}

% ============================================================
\newpage
\section{Discussion}
\label{sec:discussion}

\subsection{Supporting Future Adaptive Coordination}

% SmartWilds enables future coordination and replay research;
% do not claim the original deployment was fully closed-loop.

\subsection{Generalization Beyond Wildlife Monitoring}

% Coastal monitoring and other environmental applications.
% Keep concise and systems-focused.

\subsection{Relationship to the Ecological Analysis}

% A short boundary statement:
% The companion ecological work analyzes what the modalities reveal.
% This paper addresses how those observations were collected,
% documented, related, and made reusable.

% ============================================================
\section{Limitations}
\label{sec:limitations}

% Single site and bounded deployment.
% Some metadata reconstructed retrospectively.
% Incomplete systems telemetry.
% No claim of universal schema.
% No claim that replay substitutes for field deployment.
% Restrictions on sensitive ecological information.

% ============================================================
\section{Conclusion}
\label{sec:conclusion}




\begin{acks}
% Add field partners, ecological collaborators, data annotators,
% funding sources, and institutional support.
\end{acks}
\\

\section*{Supplemental}

% To keep the main paper from becoming a very long data manual, move these into the released artifact or appendix:

% complete dataset card;
% full field-level metadata dictionary;
% complete hardware and configuration tables;
% directory listings;
% detailed collection checklists;
% examples for every record type;
% standards crosswalk;
% extended query examples;
% full failure inventory.
\\

\bibliographystyle{ACM-Reference-Format}
\bibliography{references}

\end{document}